\chapter{Introduction, Background and Literature}
\label{ch:intro}

\section{The puzzle}
\label{sec:puzzle}

Between 2011 and 2017 the share of Indian adults holding a bank account rose from 35 to 80 per
cent \citep{findex2021}, reaching roughly 89 per cent by the 2025 Global Findex round
\citep{findex2025}. Much of that movement came from a single programme: the Pradhan Mantri Jan Dhan
Yojana (PMJDY), launched in August 2014, under which more than 56 crore accounts had been opened by
August 2025 \citep{pmjdy2025}. Measured by the indicator that dominates international league
tables, India solved financial exclusion inside a decade.

The same surveys record a second fact that sits uncomfortably beside the first. India has the
largest share of \emph{inactive} accounts of any country measured: roughly 35 per cent of Indian
adults who own an account report that it is dormant, some seven times the 5 per cent average across
other developing economies \citep{findex2021}. Official figures point the same way --- about 26 per
cent of PMJDY accounts, some 15 crore, are reported inoperative, with an average balance of roughly
\rs 4,808 \citep{pmjdy2025}. Among adults holding a dormant account, 49 per cent cite distance to
the institution and 46 per cent say they did not need an account \citep{findex2021}.

An account that exists but is never used counts as inclusion in every headline statistic while
delivering none of the consumption-smoothing, risk-pooling, or credit benefits that justify the
policy effort. The gap between \emph{having} an account and \emph{using} one is the central
empirical object of this report, and we refer to it throughout as the \textbf{access--usage gap}.

The gap could arise from low incomes, informality, distrust, poor financial literacy, or the simple
fact that many accounts were opened to receive a one-off transfer. This report examines a candidate
that policy can actually move: the infrastructure through which banking is delivered. The reasoning
is asymmetric in an important way. Opening an account is largely a \emph{supply-side} decision ---
banks can be given targets and audited against them, and Indian branch licensing has been used this
way since the 1970s. Using an account is a \emph{demand-side} decision shaped by income,
informality, and trust. Infrastructure relaxes the first constraint directly; whether it relaxes
the second is an empirical question. If the asymmetry is real, expanding infrastructure should raise
measured access more than measured usage, and the gap should \emph{widen} as infrastructure
expands. Expansion would be, in a phrase used throughout, \textbf{formal rather than functional}.

\section{Background}
\label{sec:background}

Three policy waves have shaped Indian financial inclusion, and this study's sample period sits
inside the third. Following bank nationalisation in 1969, the RBI operated a licensing rule under
which a bank could open a branch in an already-banked location only if it also opened four in
unbanked ones \citep{burgess2005}; the resulting 1977--1990 expansion remains the most studied
episode in the Indian financial-development literature. From the mid-2000s financial inclusion
became an explicit objective, bringing the no-frills account, the business correspondent model, and
the Rangarajan Committee report \citep{rangarajan2008}, together with measurement frameworks
\citep{sarma2008,sarma2012}. From 2014, PMJDY drove ownership toward saturation while the
technology stack sometimes called ``India Stack'' --- biometric identity, interoperable real-time
payments, and consented data sharing --- rewired how retail payments are made \citep{dsilva2019}.
The Unified Payments Interface processed roughly 83.7 billion transactions in calendar 2023, 131.1
billion in 2024, and 228.3 billion in 2025 \citep{npci2025}. This report's sample window sits almost
exactly on top of that transition, a point taken up in Section~\ref{sec:h1-digital}.

India also measures inclusion officially. The RBI's composite Financial Inclusion Index rose from
64.2 in March 2024 to 67.0 in March 2025 \citep{rbi_fiindex2025}, and is built from three parameters
with fixed weights: \textbf{Access 35 per cent, Usage 45 per cent, Quality 20 per cent}. Two
observations follow. First, the RBI's own scheme weights usage \emph{above} access --- an implicit
acknowledgement that account counts overstate inclusion. Second, those weights exist to collapse
three dimensions into one number, whereas this report does the opposite: it holds Access and Usage
apart precisely so their \emph{difference} can be estimated and tested. The weights are therefore
not applicable here, and Section~\ref{sec:h1-variables} records why applying them would have been
an error rather than a refinement.

\section{Literature}
\label{sec:literature}

\paragraph{Access as a supply-side lever.} The foundational result is \citet{burgess2005}, who
exploit the 1977--1990 licensing rule as policy-driven variation in rural branch openings and find
that branch expansion reduced rural poverty. The paper established both a substantive claim ---
that physical banking infrastructure has real welfare consequences --- and a methodological
template: administrative state panels, with infrastructure variation that is not simply a response
to local prosperity. The broader literature is consistent: \citet{beck2007} find financial
development disproportionately raises the incomes of the poor, and \citet{allen2016} document that
cost and proximity are among the strongest correlates of both ownership and use. One caution
prefigures our own conclusion. \citet{buliskeria2022} revisit \citet{burgess2005} and find the
poverty effect does not survive once additional policy breaks are accounted for. The lesson is not
that branches do not matter, but that in this literature the identifying variation is often thinner
and more fragile than it appears, and claims should be stress-tested against perturbations of the
sample. Chapter~\ref{ch:h1} does exactly that, and the stress test produces its most informative
result.

\paragraph{Access is not usage.} A large experimental literature establishes that removing the
barriers to \emph{opening} an account does not reliably produce \emph{use} of it. The clearest
statement is \citet{dupas2018}, who randomly offered free basic accounts in Uganda, Malawi, and
Chile: over two years only 17, 10, and 3 per cent of treated individuals respectively made five or
more deposits, and intention-to-treat effects on savings were indistinguishable from zero. Their
conclusion is the proposition this report takes to Indian administrative data --- ``policies merely
focused on expanding access to basic accounts are unlikely to improve welfare noticeably on
average.'' \citet{karlan2014} reach a compatible position, finding the binding constraints
frequently behavioural and transactional rather than a matter of account availability. This
report's contribution to the strand is one of scale and setting: the experimental evidence is
individual-level and drawn from three countries, and we ask whether the same wedge is visible in
the administrative aggregates of a country of 1.4 billion people, and whether it responds to the
infrastructure variable Indian policy actually controls.

\paragraph{Measuring inclusion.} How inclusion is measured determines what can be found.
\citet{sarma2008,sarma2012} construct a multidimensional index combining penetration, availability,
and usage; the RBI's FI-Index follows the same logic. Composite indices suit ranking and
monitoring, but they are poorly suited to the question asked here: if access and usage are averaged
into one index, the wedge between them becomes unobservable by construction --- and the wedge is
the object of interest. This report therefore departs from the index tradition deliberately.

\paragraph{Digital payments and the substitution of channels.} The final strand asks whether
physical infrastructure is still the relevant infrastructure. \citet{suri2016} show that expanded
mobile-money access in Kenya produced long-run poverty reductions operating through risk-sharing
and occupational change rather than account ownership as such. \citet{crouzet2023} show how India's
2016 demonetisation triggered durable adoption of electronic payments, shaped by network effects
rather than infrastructure stock. \citet{dsilva2019} argue that the design of India's digital public
infrastructure, rather than the density of physical bank premises, is what changed the cost of a
transaction. These matter for interpretation: if the marginal channel for financial activity
shifted from the branch to the phone during our window, a null result for branches and ATMs may be
a substantive statement about channel substitution rather than a measurement failure.

\paragraph{Where this report sits.} The experimental literature establishes the access--usage wedge
at the individual level but cannot speak to national infrastructure policy; the Indian panel
literature establishes that infrastructure has real effects but treats inclusion as a single
dimension. We combine the two --- an administrative state panel, with access and usage held apart
as separate outcomes --- so that the \emph{differential} effect of infrastructure can be estimated
directly.

\section{Hypotheses}
\label{sec:hypotheses}

The report tests three hypotheses. The first is developed and estimated in full in
Chapter~\ref{ch:h1}; the remaining two are reported in Chapters~\ref{ch:h2} and~\ref{ch:h3}.

\begin{keybox}{Hypothesis 1 (Chapter~\ref{ch:h1})}
Greater financial infrastructure raises financial inclusion, but its effect on \textbf{access} is
stronger than its effect on \textbf{usage}.

\medskip
\emph{Rationale.} Opening an account is a supply-side decision that banks can be pushed to make.
Using an account is a demand-side decision shaped by income, trust, and informality.
Infrastructure removes the first constraint, not the second.

\medskip
\emph{Expected:} $\beta_{\text{access}} > \beta_{\text{usage}} > 0$, with the difference
statistically significant.
\end{keybox}

\stub{State Hypotheses 2 and 3 in the same format --- claim, rationale, expected sign pattern ---
then develop them in Chapters~\ref{ch:h2} and~\ref{ch:h3}.}

\begin{keybox}{Hypothesis 2 (Chapter~\ref{ch:h2})}
\textit{[Statement to be supplied.]}
\end{keybox}

\begin{keybox}{Hypothesis 3 (Chapter~\ref{ch:h3})}
\textit{[Statement to be supplied.]}
\end{keybox}

\section{Contribution and roadmap}
\label{sec:contribution}

This report makes four contributions. \textbf{First}, a verified state-year panel assembled from
nine sources --- eight official series using three different year conventions, three different
state-boundary vintages and two different bank universes, plus state-level digital payment data ---
whose all-India aggregates reproduce published figures to within 0.001 per cent, with no imputation
and no missing cells. \textbf{Second}, a direct test of the access--usage wedge on administrative
data, in which the wedge is estimated as its own equation and therefore carries its own standard
error rather than being inferred by eye from two coefficients. \textbf{Third}, a decision rule
written down with its consequences before any coefficient was computed and never revised, which
proved load-bearing twice over: it delivers the verdict on physical infrastructure, and it is what
forces the honest statement that the digital result --- added after that verdict was known --- is
exploratory rather than a test that was passed. \textbf{Fourth}, a negative result that turns into
a positive one by diagnosis rather than by search. Physical infrastructure stocks move too slowly
within states for a short fixed-effects panel to identify their effect; the same diagnostic that
establishes this identifies digital payment adoption as a measure that does move, and on it the
predicted pattern appears --- running entirely through credit accounts rather than deposit
accounts, which is the reverse of what a mechanical artefact would produce.

\medskip
\noindent
Chapter~\ref{ch:data} documents the panels. Chapter~\ref{ch:h1} develops and tests Hypothesis~1 in
full, on three tracks: two measuring infrastructure physically, and one substituting digital
payment adoption. Chapters~\ref{ch:h2} and~\ref{ch:h3} present the second and third hypotheses, and
Chapter~\ref{ch:conclusion} synthesises across the three and concludes.
Appendix~\ref{app:decisions} reproduces the full decision log.
